\documentclass[14pt]{extarticle} 

% Packages for math, layout, and styling
\usepackage[utf8]{inputenc}
\usepackage{amsmath}
\usepackage{amsfonts}
\usepackage{amssymb}
\usepackage{graphicx}
\usepackage{geometry}
\usepackage{hyperref}
\usepackage{booktabs}
\usepackage{enumitem}
\usepackage{xcolor}
\usepackage{setspace}

% Page layout - standard margins with 14pt text and 1.5 spacing
\geometry{margin=1in}
\onehalfspacing 

% --- Document Information ---
\title{\textbf{Weather AI Analytics: A Lightweight Framework for Localized Forecasting via Linear Regression and Exponentially Weighted Moving Averages}}
\author{Artsiom Hancharenka, Nil Dvorakovskiy \\ \small BSU FAMCS, $3^{rd}$ group}
\date{} 

\begin{document}

\maketitle

\noindent \textbf{Abstract} \\
This paper presents a computationally efficient approach to localized meteorological prediction. By shifting away from high-resource numerical models and opaque neural networks, we propose a hybrid system utilizing Simple Linear Regression (SLR) and Exponentially Weighted Moving Averages (EWMA). The system is designed to provide explainable analytics for temperature, humidity, and precipitation. Unlike deep learning models that require vast datasets, this ``Small Data'' approach utilizes a 10-day historical window to project short-term trends. Furthermore, the inclusion of an automated validation framework allows for real-time accuracy assessment, bridging the gap between complex global models and user-centric environmental monitoring.

\bigskip

\noindent \textbf{Keywords:} \\
Explainable AI (XAI), Edge Computing, Offline Analytics, Small Data, Meteorological Heuristics, Linear Least Squares, Time-Series Forecasting.

\newpage

\noindent The modern landscape of meteorological forecasting is currently undergoing a massive shift toward Large Language Models (LLMs) and deep neural networks, yet these technologies introduce a dangerous dependency on constant internet connectivity and massive remote databases. While models like Google’s GraphCast or various GPT-based weather interfaces show impressive pattern recognition, they function essentially as ``cloud-tethered'' entities. They require a persistent, high-speed connection to reach data centers housing billions of parameters. Our research is motivated by the direct rejection of this paradigm. We believe that a weather forecasting tool is most critical precisely when the infrastructure fails—during storms, in remote agricultural zones, or in developing nations where internet stability is a luxury. Our goal is to create a product that does not rely on external databases or cloud computing, but rather functions as a standalone mathematical engine capable of running on the ``edge.''

\medskip

\noindent In an era of increasing climatic volatility, the demand for precise, localized weather information has never been higher, yet contemporary solutions frequently act as black boxes. Most consumer-grade applications present a simple icon and a percentage without revealing the age of the data or the mathematical logic behind the prediction. This lack of transparency prevents users from understanding the physical reality of their environment. By utilizing Simple Linear Regression (SLR) and Exponentially Weighted Moving Averages (EWMA), we return the power of prediction to the local device. This allows for a ``Small Data'' approach where the only requirement is a tiny rolling buffer of the last ten days of data, which can be stored in the temporary RAM of a simple microcontroller rather than a massive SQL database.

\medskip

\noindent The theoretical value of this project lies in the formalization of Persistence and Momentum Theory within a lightweight statistical framework. In short-term meteorology, we evolve the standard persistence method—which assumes tomorrow will be like today—by incorporating thermal momentum. We treat the atmosphere as a fluid system with specific heat capacity and thermal inertia. If a localized environment has shown a consistent cooling trend over several days, physical energy transfer mechanisms suggest a continuation of that vector. SLR acts as the mathematical tool to quantify this momentum, allowing us to project the most likely 24-hour state without needing to simulate the entire global atmosphere.

\medskip

\noindent While temperature changes are governed by large-scale masses, relative humidity is a high-frequency variable sensitive to immediate localized fluctuations. Humidity exhibits what we term ``short memory,'' requiring a model that prioritizes the most recent observations over distant historical data. The implementation of EWMA allows the system to prioritize recent shifts while mathematically ``forgetting'' distant history via an exponential decay function. This provides a level of reactivity that LLMs often lack, as they tend to over-smooth data based on global averages rather than reacting to a local, sudden moisture spike.
\noindent When visualizing these meteorological trends, sudden cataclysms---such
as flash floods, microbursts, or rapid cyclogenesis---manifest on the analytical
graphs as sharp, non-linear discontinuities that violate the principle of
persistence. In a standard linear least squares model, a single extreme outlier
caused by a cataclysmic event acts as a high-leverage point, disproportionately
tilting the regression slope (mm). If a localized storm produces a massive precipitation spike on the ninth day
of a ten-day window, the SLR engine interprets this not as a temporary anomaly,
but as a steepening of the overall vector. This results in a ``predictive
overshoot,'' where the model projects continued extreme conditions for the
eleventh day, even if the physical cataclysm has already dissipated. This
mathematical sensitivity highlights the inherent tension between trend-tracking
and outlier-resistance in small-data environments.
\noindent The analysis of these deviations provides critical insights into the
limitations of statistical heuristics compared to high-resource neural networks.
While an LLM might categorize a sudden deviation as sensor noise'' and smooth it
out based on global climatological averages, our EWMA-based system acknowledges
the spike as a legitimate state change. However, because EWMA relies on a
recursive decay function, a cataclysmic shift creates a recovery lag'' on the
graph. Even after the weather stabilizes, the prediction line remains anchored
to the tail-end of the cataclysm until enough new, stable data points are
ingested to overwrite the exponential memory of the event. On the user interface,
this appears as a widening gap between the projected trend and the actual
ground truth, which is immediately reflected in the Trust Score.
\noindent Ultimately, the alteration of the whole prediction by these cataclysms
serves as a diagnostic tool for the end-user. By observing a collapse in
accuracy following a sharp graphical deviation, the user is alerted to the fact
that the atmosphere has shifted from a state of momentum to a state of chaos.
This realization is vital for edge intelligence; it signals that the current
linear assumptions are no longer valid and that the user should prioritize
immediate sensor readings over the calculated 24-hour projection. In this way,
the model does not just predict the weather—it identifies its own mathematical
breaking points, providing a layer of situational awareness that remains
functional and transparent even when the most sophisticated global models are
unreachable.
\medskip
\noindent The system architecture is designed to be highly modular, ensuring that the core mathematical engine can be stripped of its web-based shell and deployed on low-power hardware. Currently, the platform uses a Flask backend to manage request orchestration and parse JSON payloads from the Open-Meteo API. However, because the forecasting logic is purely arithmetic, the backend can be entirely bypassed in a production ``Edge'' environment. In such a scenario, the system would ingest data directly from hardware sensors (thermometers and hygrometers) and process the SLR and EWMA formulas in real-time on a chip with as little as 2KB of RAM. This makes the project uniquely suited for standalone IoT devices that function in ``Offline-First'' environments.

\medskip

\noindent The integrity of the forecasting engine is currently validated through a Geocoding layer and a 10-day historical window retrieved via API. This training set provides the necessary context for the SLR and EWMA algorithms to initialize. While many AI models require years of historical data to ``train,'' our heuristic approach requires only 240 hourly data points. This efficiency is what allows the system to remain database-free; once the prediction is made, the historical data can be discarded or overwritten, keeping the storage footprint at a constant minimum. This is a radical departure from LLM-based weather services that must maintain petabytes of historical records to remain accurate.

\medskip

\noindent The mathematical framework for temperature is centered on the Least Squares Method to minimize the Residual Sum of Squares. By calculating the linear trend of temperature, we identify the slope ($m$) and intercept ($b$) of the atmospheric energy state:
\begin{equation}
m = \frac{n\sum(x_iy_i) - \sum x_i \sum y_i}{n\sum x_i^2 - (\sum x_i)^2} \quad , \quad b = \frac{\sum y_i - m \sum x_i}{n}
\end{equation}
Evaluating $y = m(n) + b$ provides the projected value for the subsequent day. This calculation is computationally trivial, requiring only basic floating-point arithmetic, which is why it can run on a digital thermometer's internal logic gate without an internet connection.

\medskip

\noindent To address the volatility of moisture and precipitation, we implement the recursive EWMA formula with a specific decay constant. By setting $\alpha = 0.7$, we ensure that 70\% of the prediction is based on the most recent 24 hours:
\begin{equation}
S_t = \alpha \cdot Y_t + (1 - \alpha) \cdot S_{t-1}
\end{equation}
This ensures the model reacts quickly to incoming weather fronts. Furthermore, we implement a physical constraint layer to prevent the mathematical model from predicting impossible states, such as negative humidity or negative rainfall:
\begin{equation}
H_{final} = \max(0, \min(100, H_{pred})) \quad , \quad P = \{0 \text{ if } P_{pred} < 0.3, \text{ else } P_{pred}\}
\end{equation}
These heuristics bridge the gap between pure mathematics and physical reality, ensuring the AI output remains within the bounds of Earth's atmosphere.

%\newpage

\noindent A primary objective of this research is to foster scientific literacy through a customized linear penalty function that acts as a Trust Score. While scientific models often use Mean Absolute Error, we employ a metric that users can intuitively understand: $A = \max(0, 100 - |T_{pred} - T_{actual}| \times 10)$. If the AI is off by $2^\circ\text{C}$, the accuracy is 80\%. If it is off by $5^\circ\text{C}$, it drops to 50\%, signaling that a non-linear event occurred—such as a sudden cold front—that the linear trend could not capture. This transparency encourages users to look for the physical causes of model failure rather than blindly trusting an algorithm.

\medskip

\noindent This framework holds significant potential for humanitarian and economic impact in developing countries where internet bandwidth is a major constraint. In many parts of the world, small-scale farmers lose crops because they cannot access localized weather forecasts due to the high cost of data or the lack of cellular coverage. By deploying this algorithm on a \$5 microcontroller connected to a simple sensor, a farmer can have a standalone weather station that predicts frost or heatwaves with high local accuracy. This is a far more sustainable solution than current trends that require a smartphone and a subscription to a cloud-based AI service.

\medskip

\noindent From an enterprise perspective, the horizontal scalability and low operational cost of this system make it an ideal microservice for private cloud environments. Organizations can process hyper-local data for thousands of specific locations—such as individual shipping containers or cell towers—without the computational overhead of global simulations. Unlike traditional forecasting services that charge per API call, our database-free engine can be embedded directly into hardware, eliminating recurring costs and data privacy concerns, as no data ever needs to leave the local network.

\medskip

\noindent The future of this framework involves the integration of barometric pressure trends to improve the detection of non-linear weather shifts. While the current linear model is highly effective for steady trends, the addition of pressure change rates will allow the system to predict the ``Turning Points'' it currently misses. By maintaining its commitment to lightweight, offline, and explainable logic, the Weather AI Analytics engine serves as a foundation for a new generation of resilient, decentralized environmental monitoring tools that prioritize user autonomy over cloud dependency.

\bigskip

\begin{thebibliography}{9}
\bibitem{openmeteo}
Zippenfenig, P. (2023). \textit{Open-Meteo: High-resolution weather API for developers}. Available at: \url{https://open-meteo.com/}.

\bibitem{regression}
Montgomery, D. C., Peck, E. A., \& Vining, G. G. (2021). \textit{Introduction to Linear Regression Analysis}. Wiley Series in Probability and Statistics.

\bibitem{ewma}
Hunter, J. S. (1986). ``The Exponentially Weighted Moving Average.'' \textit{Journal of Quality Technology}, 18(4), 203-210.
\end{thebibliography}

\end{document}